\documentclass[11pt]{article}

\usepackage[T1]{fontenc}
\usepackage{lmodern}
\usepackage[margin=1in]{geometry}
\usepackage{amsmath,amssymb,mathtools}
\usepackage{xcolor}
\usepackage{hyperref}

\setlength{\parindent}{0pt}
\setlength{\parskip}{0.5em}
\setlength{\emergencystretch}{2em}

\hypersetup{
  colorlinks=true,
  linkcolor=blue!55!black,
  urlcolor=blue!55!black,
  pdftitle={Understanding the NS blow-up-factor proof},
  pdfauthor={Machine working note}
}

\newcommand{\cH}{\mathcal H}
\newcommand{\Ibra}[1]{{}_{\mathsf I}\!\langle #1|}

\title{Understanding the NS Blow-Up-Factor Proof}
\author{Machine working note}
\date{\today}

\begin{document}

\maketitle

\section{Free-field realization}

This note follows the proof in arXiv:1312.4520, but uses standard
physicists' notation.  We denote the bosonic oscillators by \(a_n\), the NS
fermion by \(\psi_r\), and the bosonic zero modes by
\(\widehat p,\widehat q\).  Thus \(a_n\) is the oscillator denoted by
\(c_n\) in the paper.

Set
\begin{equation}
 Q=b+b^{-1},
 \qquad
 c_{\mathrm{NS}}=\frac32+3Q^2.
 \label{eq:Q-and-central-charge}
\end{equation}

\subsection{Oscillators and Fock states}

The free fields satisfy
\begin{equation}
 [a_m,a_n]=m\delta_{m+n,0},
 \qquad
 \{\psi_r,\psi_s\}=\delta_{r+s,0},
 \qquad
 [\widehat p,\widehat q]=-i,
 \label{eq:free-field-algebra}
\end{equation}
where \(m,n\in\mathbb Z\setminus\{0\}\) and
\(r,s\in\mathbb Z+\frac12\).  The momentum-\(p\) Fock vacuum obeys
\begin{equation}
 \widehat p\lvert p\rangle=p\lvert p\rangle,
 \qquad
 a_n\lvert p\rangle=0,
 \qquad
 \psi_r\lvert p\rangle=0,
 \qquad n,r>0.
 \label{eq:Fock-vacuum}
\end{equation}
The corresponding Fock space is denoted by \(\cH_p\).

\subsection{The \texorpdfstring{\(\mathsf I\)}{I}-pairing and BPZ pairing}

We distinguish the free-field Hermitian conjugation \(\mathsf I\) from
BPZ conjugation \(\dagger\).  The former is the antilinear
anti-involution defined by
\begin{equation}
 (\lambda X)^{\mathsf I}
 =\overline{\lambda}\,X^{\mathsf I},
 \qquad
 (XY)^{\mathsf I}=Y^{\mathsf I}X^{\mathsf I},
 \qquad
 \widehat p^{\mathsf I}=\widehat p,
 \qquad
 \widehat q^{\mathsf I}=\widehat q,
 \qquad
 a_n^{\mathsf I}=a_{-n},
 \qquad
 \psi_r^{\mathsf I}=\psi_{-r}.
 \label{eq:free-field-Hermitian-conjugation}
\end{equation}
In the notation of the human note,
\begin{equation}
 c_n=a_n,
 \qquad
 \eta_r=\psi_r,
 \qquad
 \widehat P=i\widehat p.
 \label{eq:human-note-free-field-dictionary}
\end{equation}
Therefore Eq.~\eqref{eq:free-field-Hermitian-conjugation} is equivalently
\begin{equation}
 c_n^{\mathsf I}=c_{-n},
 \qquad
 \eta_r^{\mathsf I}=\eta_{-r},
 \qquad
 \widehat P^{\mathsf I}=-\widehat P.
 \label{eq:I-rules-human-note-notation}
\end{equation}
The auxiliary fermion denoted by \(\psi_r\) in the human note is denoted by
\(f_r\) in this machine note.  When it is introduced below, its convention
is
\begin{equation}
 f_r^{\mathsf I}=-f_{-r},
 \qquad\text{equivalently}\qquad
 \psi_r^{\mathsf I}=-\psi_{-r}
 \quad\text{in the notation of the human note}.
 \label{eq:auxiliary-I-rule-human-note-notation}
\end{equation}
We denote the corresponding free-field Hermitian product by
\(\langle\,\cdot,\cdot\,\rangle_{\mathsf I}\), and write
\begin{equation}
 \Ibra{u}:=(|u\rangle)^{\mathsf I}.
 \label{eq:I-bra-definition}
\end{equation}
By contrast, \(\dagger\) is reserved for BPZ conjugation on an abstract
NS module:
\begin{equation}
 L_n^\dagger=L_{-n},
 \qquad
 G_r^\dagger=G_{-r}.
 \label{eq:BPZ-conjugation}
\end{equation}
For real \(Q\) and \(p\), the free-field conjugation is compatible with
the two realizations below:
\begin{equation}
 \bigl(L_n^{(\epsilon)}\bigr)^{\mathsf I}
 =L_{-n}^{(\epsilon)},
 \qquad
 \bigl(G_r^{(\epsilon)}\bigr)^{\mathsf I}
 =G_{-r}^{(\epsilon)}.
 \label{eq:I-compatible-with-SCA}
\end{equation}
Thus both pairings give the same adjoint relations for the SCA generators,
but they are defined in different places: \(\dagger\) acts on the abstract
NS module, whereas \(\mathsf I\) acts on the free-field realization.  If
\(\iota_\epsilon(p)\) is the module identification normalized by
\(\iota_\epsilon(p)|\Delta_p\rangle=|\epsilon p\rangle\), then on the real
momentum slice it may be chosen so that
\begin{equation}
 \langle u|v\rangle_{\mathrm{BPZ}}
 =
 \Ibra{\iota_\epsilon(p)u}
 \iota_\epsilon(p)v\rangle.
 \label{eq:BPZ-I-pairing-identification}
\end{equation}
The equality follows because the two sides have the same normalized
highest-weight matrix element and obey the same recursion under
\(L_n\) and \(G_r\).  Accordingly, the abstract matrix elements on the
left-hand side of the comparison below use the BPZ pairing, while their
free-field representatives on the right-hand side use the
\(\mathsf I\)-pairing.  This compatibility does not identify
\(\mathsf I\) with BPZ conjugation, especially after analytic continuation
in the momenta.

\subsection{Super-Virasoro generators}

On \(\cH_p\) there are two reflected free-field realizations, labelled by
\(\epsilon=\pm1\):
\begin{align}
 L_0^{(\epsilon)}
 &=
 \sum_{m=1}^{\infty}a_{-m}a_m
 +\sum_{r\geq\frac12}r\psi_{-r}\psi_r
 +\frac{Q^2}{8}+\frac{p^2}{2},
 \label{eq:L0-free-field}
 \\
 L_n^{(\epsilon)}
 &=
 \frac12\sum_{m\neq0,n}a_{n-m}a_m
 +\frac12\sum_{r\in\mathbb Z+\frac12}r\psi_{n-r}\psi_r
 +\left(\frac{i nQ}{2}+\epsilon p\right)a_n,
 \qquad n\neq0,
 \label{eq:Ln-free-field}
 \\
 G_r^{(\epsilon)}
 &=
 \sum_{m\neq0}a_m\psi_{r-m}
 +\left(iQr+\epsilon p\right)\psi_r.
 \label{eq:Gr-free-field}
\end{align}
They satisfy the NS super-Virasoro algebra with central charge
\eqref{eq:Q-and-central-charge}.  In either realization, \(\lvert p\rangle\)
is a highest-weight state of weight
\begin{equation}
 \Delta_p=\frac{Q^2}{8}+\frac{p^2}{2}.
 \label{eq:Delta-p}
\end{equation}

For most of the proof one chooses one realization and writes it simply as
\(L_n(p),G_r(p)\).  The second realization becomes important when the
reflection map is introduced.

\subsection{The chiral exponential vertex}

Split the chiral boson into creation and annihilation parts,
\begin{equation}
 \varphi_{<}(z)
 =-i\sum_{n=1}^{\infty}\frac{a_{-n}}{n}z^n,
 \qquad
 \varphi_{>}(z)
 =i\sum_{n=1}^{\infty}\frac{a_n}{n}z^{-n},
 \label{eq:boson-splitting}
\end{equation}
and define
\begin{equation}
 \psi(z)
 =\sum_{r\in\mathbb Z+\frac12}\psi_r z^{-r-\frac12}.
 \label{eq:fermion-field}
\end{equation}

The free-field vertex operator of momentum \(\alpha\) is the ordered
exponential
\begin{equation}
 \boxed{
 E^{\alpha}(z)
 =z^{-\Delta_\alpha}
 e^{\frac{\alpha}{2}\widehat q}
 e^{\alpha\varphi_{<}(z)}
 z^{-i\alpha\widehat p}
 e^{\alpha\varphi_{>}(z)}
 e^{\frac{\alpha}{2}\widehat q}
 },
 \qquad
 \Delta_\alpha=\frac12\alpha(Q-\alpha).
 \label{eq:ordered-exponential}
\end{equation}
In physicists' shorthand,
\begin{equation}
 E^{\alpha}(z)
 =z^{-\Delta_\alpha}:e^{\alpha\varphi(z)}:,
 \label{eq:normal-ordered-shorthand}
\end{equation}
where the normal-ordering symbol includes the symmetric zero-mode ordering
displayed explicitly in Eq.~\eqref{eq:ordered-exponential}.

The two basic checks are
\begin{align}
 [\widehat p,E^{\alpha}(z)]
 &=-i\alpha E^{\alpha}(z),
 \label{eq:E-momentum-shift}
 \\
 [L_n(p),E^{\alpha}(z)]
 &=z^n\left(z\partial_z+(n+1)\Delta_\alpha\right)E^{\alpha}(z),
 \label{eq:E-Virasoro-primary}
 \\
 [G_r(p),E^{\alpha}(z)]
 &=-i\alpha z^{r+\frac12}\psi(z)E^{\alpha}(z).
 \label{eq:E-superpartner}
\end{align}
Hence \(E^{\alpha}(z)\) is the free-field representative of an NS primary
vertex and shifts the Fock momentum according to
\begin{equation}
 E^{\alpha}(z):\cH_p\longrightarrow\cH_{p-i\alpha}.
 \label{eq:E-map}
\end{equation}

At this stage no screening charge has been introduced.  The ordered
exponential is the local chiral primary vertex; screening charges will enter
only when we ask for a nonzero matrix element between Fock spaces whose
momenta do not satisfy the neutrality condition for \(E^{\alpha}\) alone.

\section{Why the screened matrix element represents an NS primary}

The goal of this section is to explain the identity written as Eq.~(2.9) in
arXiv:1312.4520.  The logic has three steps:
\begin{enumerate}
 \item define a screened matrix element in the free-field Fock space and
       determine its SCA properties;
 \item define the corresponding matrix element of an abstract NS primary
       chiral vertex;
 \item compare the two constructions and fix the normalization for which
       they agree exactly.
\end{enumerate}

\subsection{Screening charges}

For either
\begin{equation}
 \beta=b
 \qquad\text{or}\qquad
 \beta=b^{-1},
 \label{eq:screening-momenta}
\end{equation}
one has
\begin{equation}
 \Delta_\beta
 =\frac12\beta(Q-\beta)
 =\frac12.
 \label{eq:screening-weight-half}
\end{equation}
Consequently,
\begin{equation}
 S_\beta(z):=\psi(z)E^\beta(z)
 \label{eq:screening-current}
\end{equation}
is a fermion-odd field of conformal weight one.  Its transformations are
total derivatives:
\begin{align}
 [L_n(p),S_\beta(z)]
 &=
 \partial_z\!\left(z^{n+1}S_\beta(z)\right),
 \label{eq:screening-L-total-derivative}
 \\
 \{G_r(p),S_\beta(z)\}
 &=
 \frac{i}{\beta}
 \partial_z\!\left(z^{r+\frac12}E^\beta(z)\right).
 \label{eq:screening-G-total-derivative}
\end{align}
For a closed Coulomb-gas contour, define
\begin{equation}
 \mathcal Q_\beta
 :=
 \oint dz\,S_\beta(z)
 =
 \oint dz\,\psi(z)E^\beta(z).
 \label{eq:screening-charge-general}
\end{equation}
The total derivatives imply
\begin{equation}
 [L_n(p),\mathcal Q_\beta]=0,
 \qquad
 \{G_r(p),\mathcal Q_\beta\}=0.
 \label{eq:screening-supercommutes}
\end{equation}
Thus a screening charge shifts the Fock momentum while supercommuting with
the SCA:
\begin{equation}
 \mathcal Q_b:\mathcal H_p\longrightarrow\mathcal H_{p-ib},
 \qquad
 \mathcal Q_{1/b}:\mathcal H_p\longrightarrow\mathcal H_{p-i/b}.
 \label{eq:screening-momentum-shifts}
\end{equation}

\subsection{Step 1: the screened free-field matrix element}

Fix nonnegative integers \(r,s\).  For the even NS primary component, assume
\begin{equation}
 r+s\in2\mathbb Z,
 \label{eq:even-screening-number}
\end{equation}
so that the product of screening charges is even.  Define the screened
free-field operator
\begin{equation}
 \mathcal I_{\alpha;r,s}(z)
 :=
 E^\alpha(z)\mathcal Q_b^r\mathcal Q_{1/b}^s
 \label{eq:screened-free-field-operator}
\end{equation}
and its matrix element
\begin{equation}
 \mathcal F_{\alpha;r,s}^{\pm,\pm}(q,p;z)
 :=
 \Ibra{\pm q}
 E^\alpha(z)\mathcal Q_b^r\mathcal Q_{1/b}^s
 |\pm p\rangle.
 \label{eq:screened-free-field-matrix-element}
\end{equation}
The signs in front of \(p\) and \(q\) are independent.

The operator in Eq.~\eqref{eq:screened-free-field-operator} shifts the
momentum by
\begin{equation}
 -i\left(\alpha+rb+sb^{-1}\right).
 \label{eq:first-total-momentum-shift}
\end{equation}
Therefore Eq.~\eqref{eq:screened-free-field-matrix-element} can be nonzero
only on the neutrality locus
\begin{equation}
 \boxed{
 \pm q
 =
 \pm p-i\alpha-i\left(rb+sb^{-1}\right)
 }.
 \label{eq:first-neutrality-condition}
\end{equation}

The screening charges do not alter the SCA transformation law of
\(E^\alpha\).  Hence
\begin{equation}
 [L_n(p),\mathcal I_{\alpha;r,s}(z)]
 =
 z^n\left(
 z\partial_z+(n+1)\Delta_\alpha
 \right)\mathcal I_{\alpha;r,s}(z),
 \label{eq:screened-L-transformation}
\end{equation}
where
\begin{equation}
 \Delta_\alpha=\frac12\alpha(Q-\alpha).
 \label{eq:abstract-primary-weight}
\end{equation}
Define the odd partner
\begin{equation}
 \mathcal I_{\alpha;r,s}^{*}(z)
 :=
 -i\alpha\,
 \psi(z)E^\alpha(z)
 \mathcal Q_b^r\mathcal Q_{1/b}^s.
 \label{eq:screened-odd-partner}
\end{equation}
Then
\begin{align}
 [G_r(p),\mathcal I_{\alpha;r,s}(z)]
 &=
 z^{r+\frac12}\mathcal I_{\alpha;r,s}^{*}(z),
 \label{eq:screened-G-transformation}
 \\
 [L_n(p),\mathcal I_{\alpha;r,s}^{*}(z)]
 &=
 z^n\left(
 z\partial_z+(n+1)\left(\Delta_\alpha+\frac12\right)
 \right)\mathcal I_{\alpha;r,s}^{*}(z),
 \label{eq:screened-star-L-transformation}
 \\
 \{G_r(p),\mathcal I_{\alpha;r,s}^{*}(z)\}
 &=
 z^{r-\frac12}
 \left(
 z\partial_z+(2r+1)\Delta_\alpha
 \right)\mathcal I_{\alpha;r,s}(z).
 \label{eq:screened-star-G-transformation}
\end{align}
These are precisely the transformation laws of an NS primary
supermultiplet.

In particular, the \(L_0\) Ward identity fixes the coordinate dependence:
\begin{equation}
 \mathcal F_{\alpha;r,s}^{\pm,\pm}(q,p;z)
 =
 z^{\Delta_q-\Delta_\alpha-\Delta_p}
 \mathcal F_{\alpha;r,s}^{\pm,\pm}(q,p;1),
 \label{eq:free-field-primary-z-dependence}
\end{equation}
where
\begin{equation}
 \Delta_p=\frac{Q^2}{8}+\frac{p^2}{2},
 \qquad
 \Delta_q=\frac{Q^2}{8}+\frac{q^2}{2}.
 \label{eq:p-and-q-weights}
\end{equation}

\subsection{Step 2: the abstract NS-primary matrix element}

Let \(\mathcal V_{\Delta_p}\) and \(\mathcal V_{\Delta_q}\) be NS
super-Virasoro highest-weight modules.  Their highest-weight states satisfy
\begin{align}
 L_0|\Delta_p\rangle
 &=\Delta_p|\Delta_p\rangle,
 &
 L_n|\Delta_p\rangle
 &=
 G_r|\Delta_p\rangle=0,
 \qquad n,r>0,
 \label{eq:incoming-NS-primary}
 \\
 \langle\Delta_q|L_0
 &=\Delta_q\langle\Delta_q|,
 &
 \langle\Delta_q|L_{-n}
 &=
 \langle\Delta_q|G_{-r}=0,
 \qquad n,r>0.
 \label{eq:outgoing-NS-primary}
\end{align}
We use unit normalization in each fixed module:
\begin{equation}
 \langle\Delta_p|\Delta_p\rangle
 =
 \langle\Delta_q|\Delta_q\rangle
 =
 \langle\Delta_\alpha|\Delta_\alpha\rangle
 =1.
 \label{eq:abstract-highest-state-normalization}
\end{equation}

An even NS primary chiral vertex \(V_{\Delta_\alpha}(z)\), together with its
odd superpartner \(V_{\Delta_\alpha}^{*}(z)\), is defined by
\begin{align}
 [L_n,V_{\Delta_\alpha}(z)]
 &=
 z^n\left(
 z\partial_z+(n+1)\Delta_\alpha
 \right)V_{\Delta_\alpha}(z),
 \label{eq:abstract-primary-L-transformation}
 \\
 [G_r,V_{\Delta_\alpha}(z)]
 &=
 z^{r+\frac12}V_{\Delta_\alpha}^{*}(z),
 \label{eq:abstract-primary-G-transformation}
 \\
 [L_n,V_{\Delta_\alpha}^{*}(z)]
 &=
 z^n\left(
 z\partial_z+(n+1)\left(\Delta_\alpha+\frac12\right)
 \right)V_{\Delta_\alpha}^{*}(z),
 \label{eq:abstract-star-L-transformation}
 \\
 \{G_r,V_{\Delta_\alpha}^{*}(z)\}
 &=
 z^{r-\frac12}
 \left(
 z\partial_z+(2r+1)\Delta_\alpha
 \right)V_{\Delta_\alpha}(z).
 \label{eq:abstract-star-G-transformation}
\end{align}
The primary matrix element is
\begin{equation}
 \mathcal A_\alpha(q,p;z)
 :=
 \langle\Delta_q|
 V_{\Delta_\alpha}(z)
 |\Delta_p\rangle.
 \label{eq:abstract-primary-matrix-element}
\end{equation}
Its \(L_0\) Ward identity gives
\begin{equation}
 \mathcal A_\alpha(q,p;z)
 =
 z^{\Delta_q-\Delta_\alpha-\Delta_p}
 \mathcal A_\alpha(q,p;1).
 \label{eq:abstract-primary-z-dependence}
\end{equation}
The remaining SCA Ward identities determine every descendant matrix element
in terms of the single constant
\begin{equation}
 C_{q\alpha p}
 :=
 \langle\Delta_q|
 V_{\Delta_\alpha}(1)
 |\Delta_p\rangle.
 \label{eq:abstract-chiral-vertex-constant}
\end{equation}
The unit norms of the three primary states do not determine
\(C_{q\alpha p}\); it is the remaining normalization of the chiral vertex.

\subsection{Step 3: comparison and exact normalization}

Normalize the Fock vacua in each fixed-momentum fiber by
\begin{equation}
 \Ibra{\pm p}\pm p\rangle
 =
 \Ibra{\pm q}\pm q\rangle
 =1.
 \label{eq:Fock-vacuum-normalization}
\end{equation}
Before restricting to fixed fibers, this is the usual delta-function
normalization
\begin{equation}
 \Ibra{p_1}p_2\rangle=\delta(p_1-p_2).
 \label{eq:delta-function-normalization}
\end{equation}
Fix the module identifications by
\begin{equation}
 \iota_\pm(p)|\Delta_p\rangle=|\pm p\rangle,
 \qquad
 \iota_\pm(q)|\Delta_q\rangle=|\pm q\rangle.
 \label{eq:normalized-module-identifications}
\end{equation}
For generic real momenta these extend to unitary isomorphisms between the
abstract NS Verma modules and their free-field realizations.

The screened operator \(\mathcal I_{\alpha;r,s}\) and the abstract chiral
vertex \(V_{\Delta_\alpha}\) now have
\begin{enumerate}
 \item the same NS superconformal transformation laws,
 \item the same conformal weight \(\Delta_\alpha\),
 \item the same fermion parity, and
 \item compatible incoming and outgoing modules on the neutrality locus
       \eqref{eq:first-neutrality-condition}.
\end{enumerate}
They therefore define the same chiral intertwiner up to one overall
constant.  Exact equality is obtained by choosing
\begin{equation}
 \boxed{
 C_{q\alpha p}
 =
 \mathcal F_{\alpha;r,s}^{\pm,\pm}(q,p;1)
 =
 \Ibra{\pm q}
 E^\alpha(1)\mathcal Q_b^r\mathcal Q_{1/b}^s
 |\pm p\rangle
 }.
 \label{eq:compatible-chiral-vertex-normalization}
\end{equation}
With this normalization,
\begin{equation}
 \boxed{
 \langle\Delta_q|
 V_{\Delta_\alpha}(z)
 |\Delta_p\rangle
 =
 \Ibra{\pm q}
 E^\alpha(z)\mathcal Q_b^r\mathcal Q_{1/b}^s
 |\pm p\rangle
 },
 \label{eq:first-form-of-2.9}
\end{equation}
which is the first line of Eq.~(2.9).

For descendant states \(u\in\mathcal V_{\Delta_q}\) and
\(v\in\mathcal V_{\Delta_p}\), the normalization-independent form of the
same statement is
\begin{equation}
 \boxed{
 \frac{
  \langle u|V_{\Delta_\alpha}(1)|v\rangle
 }{
  \langle\Delta_q|V_{\Delta_\alpha}(1)|\Delta_p\rangle
 }
 =
 \frac{
  \Ibra{\iota_\pm(q)u}
  E^\alpha(1)\mathcal Q_b^r\mathcal Q_{1/b}^s
  |\iota_\pm(p)v\rangle
 }{
  \Ibra{\pm q}
  E^\alpha(1)\mathcal Q_b^r\mathcal Q_{1/b}^s
  |\pm p\rangle
 }
 }.
 \label{eq:normalized-screened-Ward-identity}
\end{equation}
Thus the unit normalization of the states is necessary, but it is not
sufficient by itself: one must also fix the remaining chiral-vertex constant
as in Eq.~\eqref{eq:compatible-chiral-vertex-normalization}.

\subsection{The reflected exponential and the eight representations}

Since
\begin{equation}
 \Delta_{Q-\alpha}
 =
 \frac12(Q-\alpha)\alpha
 =
 \Delta_\alpha,
 \label{eq:reflected-alpha-weight}
\end{equation}
the exponential \(E^{Q-\alpha}\) has the same SCA primary weight.  Its total
momentum shift in the presence of the same screening charges is
\begin{equation}
 -i\left(Q-\alpha+rb+sb^{-1}\right)
 =
 i\alpha-i\left((r+1)b+(s+1)b^{-1}\right).
 \label{eq:second-total-momentum-shift}
\end{equation}
The second neutrality locus is therefore
\begin{equation}
 \boxed{
 \pm q
 =
 \pm p+i\alpha
 -i\left((r+1)b+(s+1)b^{-1}\right)
 }.
 \label{eq:second-neutrality-condition}
\end{equation}
With the analogous choice of chiral-vertex normalization,
\begin{equation}
 \boxed{
 \langle\Delta_q|
 V_{\Delta_\alpha}(z)
 |\Delta_p\rangle
 =
 \Ibra{\pm q}
 E^{Q-\alpha}(z)\mathcal Q_b^r\mathcal Q_{1/b}^s
 |\pm p\rangle
 }.
 \label{eq:second-form-of-2.9}
\end{equation}

The abstract weights depend only on \(p^2\) and \(q^2\), so the signs in
front of \(p\) and \(q\) can be chosen independently.  There are four sign
choices and two exponential choices,
\begin{equation}
 E^\alpha
 \qquad\text{or}\qquad
 E^{Q-\alpha}.
\end{equation}
This gives the eight free-field representations stated below Eq.~(2.9).

\subsection{The precise meaning of Eq.~(2.9)}

Equation~(2.9) is a statement about matrix elements on the charge-neutral
loci \eqref{eq:first-neutrality-condition} and
\eqref{eq:second-neutrality-condition}.  It is not one global operator
identity
\begin{equation}
 V_{\Delta_\alpha}(z)
 =
 E^\alpha(z)\mathcal Q_b^r\mathcal Q_{1/b}^s
 \label{eq:not-global-operator-identity}
\end{equation}
valid for fixed \(r,s\) on every Fock module.  Rather, each admissible
screened expression realizes the same normalized chiral intertwiner on its
own neutrality locus.  Later, the blow-up-factor proof evaluates matrix
elements on many such loci and uses polynomiality to continue the result to
generic momenta.

For the even primary matrix element in Eq.~(2.9), the total number \(r+s\)
of screening charges is even.  For the odd superpartner matrix elements in
Eqs.~(2.10) and (2.11), it is odd.

\section{From the screened intertwiner to Eq.~(3.13)}

We now explain the precise chain of substitutions that turns the normalized
three-point block into the free-field matrix element in Eq.~(3.13).  The
essential input is not an explicit formula for the reflected fermion.
Instead, one first works in a free-field chart in which no reflected
fermions occur.

\subsection{The unreflected \texorpdfstring{\(\chi\)}{chi}-states}

Introduce the additional fermion \(f_r\) by
\begin{equation}
 \{f_r,f_s\}=\delta_{r+s,0},
 \qquad
 f_r^{\mathsf I}=-f_{-r},
 \label{eq:auxiliary-fermion-I-rule}
\end{equation}
and define
\begin{equation}
 \chi_r=f_r-i\psi_r.
 \label{eq:chi-mode-definition}
\end{equation}
The two fermions anticommute with one another.  It follows that
\begin{equation}
 \chi_r^{\mathsf I}=-\chi_{-r},
 \qquad
 \{\chi_r,\chi_s\}=0.
 \label{eq:chi-I-and-anticommutator}
\end{equation}
Let
\begin{equation}
 |p,0\rangle:=|p\rangle\otimes|0_f\rangle.
 \label{eq:extended-Fock-vacuum}
\end{equation}
For \(j>0\), the double-Virasoro highest-weight state has the unreflected
free-field representative
\begin{equation}
 |p,j\rangle
 =
 \chi_{-j+\frac12}\cdots
 \chi_{-\frac32}\chi_{-\frac12}|p,0\rangle.
 \label{eq:positive-j-unreflected-state}
\end{equation}
The pole-cleared state used in the three-point block is
\begin{equation}
 |\xi_{p,j}\rangle
 :=
 \Omega(p,j)|p,j\rangle,
 \label{eq:xi-state-definition}
\end{equation}
where \(\Omega(p,j)\) is the polynomial introduced in Eq.~(2.24) of the
paper.  For real \(p,Q,b\),
\begin{equation}
 \Omega(p,j)^{\mathsf I}
 =
 \overline{\Omega(p,j)}
 =
 \Omega(-p,j).
 \label{eq:Omega-I-conjugation}
\end{equation}
Consequently, the free-field Hermitian bra associated with
Eq.~\eqref{eq:xi-state-definition} is
\begin{equation}
 \bigl(|\xi_{p,j}\rangle\bigr)^{\mathsf I}
 =
 \Omega(-p,j)\,
 \Ibra{p,0}
 \chi_{-\frac12}^{\mathsf I}
 \chi_{-\frac32}^{\mathsf I}\cdots
 \chi_{-j+\frac12}^{\mathsf I}.
 \label{eq:positive-j-I-bra}
\end{equation}
The replacement \(p\mapsto-p\) in the coefficient in
Eq.~\eqref{eq:positive-j-I-bra} is caused by complex conjugation under
\(\mathsf I\).  It is not an application of the reflection map.

\subsection{Turning the middle state into contour insertions}

The local field
\begin{equation}
 \chi(\xi)
 =
 \sum_{r\in\mathbb Z+\frac12}
 \chi_r\,\xi^{-r-\frac12}
 \label{eq:chi-local-field}
\end{equation}
implements the \(\chi\)-modes through the state-operator correspondence.
Around an insertion at \(z=1\),
\begin{equation}
 \chi_{-\ell+\frac12}\mathcal O(1)
 =
 \frac{1}{2\pi i}
 \oint_{1}
 \frac{d\xi}{(\xi-1)^\ell}
 \chi(\xi)\mathcal O(1),
 \qquad
 \ell=1,2,\ldots.
 \label{eq:chi-mode-contour-extraction}
\end{equation}
Therefore the field associated with the middle state can be written as
\begin{align}
 V_{\xi_{p_2,j_2}}(1)
 ={}&
 \frac{\Omega(p_2,j_2)}{(2\pi i)^{j_2}}
 \oint_{1}\frac{d\xi_{j_2}}{(\xi_{j_2}-1)^{j_2}}
 \cdots
 \oint_{1}\frac{d\xi_1}{\xi_1-1}
 \nonumber\\
 &\quad\times
 \chi(\xi_{j_2})\cdots\chi(\xi_1)
 V_{p_2,0}(1).
 \label{eq:middle-chi-state-operator-map}
\end{align}
The contours and the fermionic fields are ordered as displayed.  This
ordering is inherited from the ordered product of modes in the state and
must be kept fixed when signs are evaluated.

\subsection{Replacing the primary vertex by a screened intertwiner}

Choose a screening realization satisfying
\begin{equation}
 p_3
 =
 p_1-i\left(\alpha_2+rb+sb^{-1}\right),
 \label{eq:3.13-neutrality-condition}
\end{equation}
and abbreviate
\begin{equation}
 \mathcal S_{\alpha_2;r,s}
 :=
 E^{\alpha_2}(1)\mathcal Q_b^r\mathcal Q_{1/b}^s,
 \qquad
 \mathcal N_{31}
 :=
 \Ibra{p_3}
 \mathcal S_{\alpha_2;r,s}
 |p_1\rangle.
 \label{eq:screened-intertwiner-and-normalization}
\end{equation}
The Ward-identity argument of the previous section identifies the
normalized abstract primary vertex with the normalized screened
intertwiner:
\begin{equation}
 \frac{
  V_{p_2,0}(1)
 }{
  \langle p_3|V_{p_2,0}(1)|p_1\rangle
 }
 \quad\longleftrightarrow\quad
 \frac{
  \mathcal S_{\alpha_2;r,s}
 }{
  \mathcal N_{31}
 }.
 \label{eq:normalized-primary-to-screened-intertwiner}
\end{equation}
This is an equality of normalized chiral intertwiners on the neutrality
locus, not a global operator identity.  Once the primary matrix element is
fixed, the SCA Ward identities guarantee equality of all descendant matrix
elements.

\subsection{Substitution into the normalized three-point block}

Suppose \(j_1,j_2,j_3>0\) and
\(j_1+j_2+j_3\in2\mathbb Z\).  The normalized block is
\begin{equation}
 \rho_{\mathrm{NS}}^{\mathcal A}
 \left(
  \xi_{p_3,j_3},
  \xi_{p_2,j_2},
  \xi_{p_1,j_1}
  \mathrel{|}1
 \right)
 :=
 \frac{
  \langle\xi_{p_3,j_3}|
  V_{\xi_{p_2,j_2}}(1)
  |\xi_{p_1,j_1}\rangle
 }{
  \langle p_3|V_{p_2,0}(1)|p_1\rangle
 }.
 \label{eq:normalized-A-three-point-block}
\end{equation}
The bras in Eq.~\eqref{eq:normalized-A-three-point-block} are abstract BPZ
bras.  After applying the normalized module identification, the
right-hand free-field calculation uses the \(\mathsf I\)-pairing.

Substituting Eqs.~\eqref{eq:positive-j-unreflected-state},
\eqref{eq:positive-j-I-bra},
\eqref{eq:middle-chi-state-operator-map}, and
\eqref{eq:normalized-primary-to-screened-intertwiner} gives
\begin{align}
 &\rho_{\mathrm{NS}}^{\mathcal A}
 \left(
  \xi_{p_3,j_3},
  \xi_{p_2,j_2},
  \xi_{p_1,j_1}
  \mathrel{|}1
 \right)
 \nonumber\\
 &=
 \frac{
  \Omega(-p_3,j_3)\Omega(p_2,j_2)\Omega(p_1,j_1)
 }{
  (2\pi i)^{j_2}\mathcal N_{31}
 }
 \oint_{1}\frac{d\xi_{j_2}}{(\xi_{j_2}-1)^{j_2}}
 \cdots
 \oint_{1}\frac{d\xi_1}{\xi_1-1}
 \nonumber\\
 &\quad\times
 \Ibra{p_3,0}
 \chi_{-\frac12}^{\mathsf I}\cdots
 \chi_{-j_3+\frac12}^{\mathsf I}
 \chi(\xi_{j_2})\cdots\chi(\xi_1)
 \mathcal S_{\alpha_2;r,s}
 \nonumber\\
 &\hspace{8em}\times
 \chi_{-j_1+\frac12}\cdots\chi_{-\frac12}
 |p_1,0\rangle.
 \label{eq:3.13-with-I-conjugation}
\end{align}
This is Eq.~(3.13), with the paper's free-field dagger replaced by our
\(\mathsf I\).  Expanding
\(\chi_r^{\mathsf I}=-\chi_{-r}\) produces the factor
\((-1)^{j_3}\) that appears when the correlator is separated into its
bosonic and fermionic parts.

\subsection{What happened to the reflected fermion?}

Equation~\eqref{eq:3.13-with-I-conjugation} does not express
\(\psi_r^R(p)\) in terms of the original fermions.  It avoids that problem:
\begin{enumerate}
 \item For \(j_1,j_2,j_3>0\), all three states are represented in the
       unreflected \(\chi\)-basis.
 \item Other sign choices are analyzed using the alternative free-field
       realizations associated with \(E^{\alpha}\), \(E^{Q-\alpha}\), and
       the independent signs of the incoming and outgoing momenta.
 \item The factors \(\Omega(p_i,j_i)\) remove the possible poles and turn
       the normalized block into a polynomial of bounded degree.
 \item Free-field calculations on the admissible screening loci determine
       the zeros and normalization of this polynomial, after which the
       result is continued to generic momenta.
\end{enumerate}
Thus the authors circumvent the absence of a closed formula for the
reflected fermion; they do not solve for it in the original oscillator
basis.  The calculation succeeds because each required specialization can
be evaluated in a free-field chart adapted to that specialization.

\end{document}
